\pdfoutput=1
% nopdfoutputerror: quantumarticle v6.1's pdfoutput check mis-fires on the
% TeX Live 2023 kernel even with \pdfoutput=1 set; we compile with real
% pdflatex, so the guard is unneeded. Remove if the class is updated.
\documentclass[a4paper,onecolumn,unpublished,nopdfoutputerror,allowtoday]{quantumarticle}

% ---------------------------------------------------------------------
%  Venue-specialized port for the journal Quantum (quantum-journal.org).
%  GENERATED from the base ../../main.tex (the venue-neutral source of
%  truth); keep the math and body in sync with the base. quantumarticle
%  supplies geometry, fonts, hyperref, and amsmath.
% ---------------------------------------------------------------------

\usepackage{amssymb}
\usepackage{amsthm}
\usepackage{mathtools}
\usepackage{booktabs}
\usepackage{array}
\usepackage{enumitem}
\usepackage[numbers,square]{natbib}
\bibliographystyle{plainnat}
\usepackage{hyperref}   % load explicitly so cleveref orders after it (quantumarticle also loads hyperref; second load is a no-op)
\usepackage{cleveref}

\numberwithin{equation}{section}

% --- theorem environments ---
\theoremstyle{plain}
\newtheorem{theorem}{Theorem}[section]
\newtheorem{proposition}[theorem]{Proposition}
\newtheorem{lemma}[theorem]{Lemma}
\newtheorem{corollary}[theorem]{Corollary}
\theoremstyle{definition}
\newtheorem{definition}[theorem]{Definition}
\newtheorem{example}[theorem]{Example}
\theoremstyle{remark}
\newtheorem{remark}[theorem]{Remark}

% --- notation ---
\newcommand{\fpr}{\epsilon}          % false-positive rate
\newcommand{\fnr}{\omega}            % false-negative rate
\newcommand{\eps}{\epsilon}
\newcommand{\Prob}{\mathbb{P}}
\newcommand{\Ind}[1]{\mathbf{1}\!\left[#1\right]}
\newcommand{\given}{\,\vert\,}
\newcommand{\ket}[1]{\lvert #1 \rangle}
\newcommand{\bra}[1]{\langle #1 \rvert}
\newcommand{\bits}{\{0,1\}}
\DeclareMathOperator{\diag}{diag}
\DeclareMathOperator{\tr}{tr}
\newcommand{\bxor}{\oplus}
\newcommand{\bigxor}{\bigoplus}
\newcommand{\tensor}{\otimes}

\title{A Forward \texorpdfstring{$(\eps,\omega)$}{(epsilon, omega)} Calculus for
Quantum Readout Error and the Classical Post-Processing of Measurement Outcomes}

\author{Alexander Towell}
\affiliation{Southern Illinois University Edwardsville}
\email{atowell@siue.edu}
\orcid{0000-0001-6443-9897}

\date{\today}

\begin{document}
\maketitle

\begin{abstract}
Quantum readout (measurement) error is, on the diagonal sector of the density
matrix, a classical bit-flip channel: a measured qubit is a truth value
garbled by an asymmetric channel with a false-positive rate $\eps$ and a
false-negative rate $\omega$. The two rates genuinely differ, because $T_1$
relaxation during the measurement window makes $\ket{1}\!\to\!\ket{0}$ misreads
more likely than the reverse, so $\omega > \eps$. Yet the bits a circuit
\emph{reports} are rarely raw measurements; they are Boolean functions of
measured ancillas, parities for syndrome and stabilizer extraction, majority
votes, and feed-forward conditionals, and the error on those derived bits is
today reconstructed numerically per circuit, or estimated with a symmetrized,
single-rate rule of thumb. We give the missing piece: a reusable,
\emph{forward}, \emph{asymmetric} calculus that propagates a pair $(\eps,\omega)$
in closed form through Boolean combinations (\textsc{not}, \textsc{and},
\textsc{or}, parity) and through independent-instance (Kronecker) composition of
qubits. The calculus is exactly the error model of the Bernoulli random
approximate set and map framework, so it also unifies a quantum-hardware error
model with the classical theory of Bloom filters and cipher maps. Our lead
example is the reliability of a quantum-error-correction \emph{syndrome bit}: a
parity of $n$ noisy ancilla measurements has flip probability exactly
$\tfrac12\bigl(1 - \prod_{i=1}^{n}(1-2q_i)\bigr)$, which we contrast with the
field's casual ``about $n$ times the per-measurement error'' linearization.
The linear rule already overstates the true error by up to tens of percent at
modest $n$ and per-measurement rates. The contribution is \emph{forward} prediction
and error budgeting, complementary to inverse mitigation such as M3; the closed
forms hold under per-bit channel independence, with correlated-channel
(crosstalk) extensions left as future work. We position the calculus carefully
against the matrix-inversion mitigation paradigm, effectus theory, the classical
noisy-Boolean-formula identities, asymmetric quantum hypothesis testing, and the
relaxation-physics origin of the readout asymmetry.
\end{abstract}

\section{Introduction: the gap}
\label{sec:intro}

A near-term quantum processor does not hand the programmer the bits it computed.
It hands the programmer \emph{noisy readouts} of those bits, and then, in any
nontrivial program, the classical control system computes \emph{further}
Boolean functions of those noisy readouts before anything is reported or acted
upon. Error-correction cycles extract \emph{syndrome bits}, each the parity of
several measured ancilla qubits. Repetition and majority schemes take the
\emph{majority vote} of repeated measurements. Adaptive circuits feed a measured
bit \emph{forward} into a classically controlled gate. In every case the bit
that matters is a function $g(b_1,\dots,b_n)$ of noisy measured bits, and the
quantity of interest is the error rate of $g$.

Two facts about readout error make this propagation subtle, and the literature
handles each only partially.

\paragraph{Readout error is asymmetric.} The probability of reading $0$ when the
qubit is really $1$ is not equal to the probability of reading $1$ when it is
really $0$. The dominant mechanism is $T_1$ (energy) relaxation during the finite
measurement window: an excited $\ket{1}$ can decay to $\ket{0}$ mid-measurement,
biasing errors toward $1\!\to\!0$. Measured readout error on superconducting
hardware runs from a few percent to tens of percent and is reliably
state-dependent \citep{tannu2019mitigating,hicks2021readout}. A correct error
model for a derived bit must therefore carry \emph{two} rates, not one, and must
carry them through the Boolean post-processing without collapsing them into a
single average. Much of the applied tooling nonetheless symmetrizes, reporting a
single per-qubit error and losing the very asymmetry that the physics produces.

\paragraph{Derived-bit error is reconstructed numerically.} The mature response
to readout error is \emph{mitigation}: build an assignment (confusion) matrix $A$
relating ideal to observed outcome distributions and (approximately) invert it.
This is the matrix-inversion paradigm, with scalable realizations such as M3
\citep{nation2021m3} and the multiqubit treatment of \citet{bravyi2021mitigating}.
Mitigation is an \emph{inverse} procedure: given an observed distribution, recover
an estimate of the ideal one. It is numerical, it is per-circuit, and it answers
a different question from the one a circuit designer asks at design time, namely:
\emph{given} per-qubit readout rates, what is the predicted error on this
particular derived bit, symbolically and before any data are taken? There is at
present no reusable, closed-form, forward calculus that answers this question
while preserving the asymmetry.

\paragraph{Contribution.} We supply that calculus. Modeling a measured bit as a
$2\times2$ row-stochastic channel with false-positive rate $\eps$ and
false-negative rate $\omega$, we derive closed forms for how the pair
$(\eps,\omega)$ propagates through the Boolean connectives and through
independent (Kronecker) composition of qubits, carrying both rates
asymmetrically through an entire Boolean tree. The calculus is precisely the
error model of \emph{Bernoulli type theory}, the theory of random approximate
sets, maps, and relations that has Bloom filters, perfect hash filters, and
cipher maps as special cases \citep{bernoulliSets,bernoulliComposition,bernoulliMaps}.
The single-bit asymmetric channel and the fact that independent-qubit errors
tensor are \emph{not} new (see \cref{sec:related}); the new artifact is the
\emph{forward asymmetric Boolean calculus}, its application to quantum readout,
and the unification it exposes between a quantum-hardware error model and a
classical probabilistic-data-structure error model.

\paragraph{Why this is sound, and where it stops.} The model treats the qubit
as carrying a definite latent classical value that the channel garbles. That is
faithful exactly on the \emph{diagonal} sector of the density matrix, the
readout sector, where the state has no off-diagonal coherences to interfere.
A latent-value-plus-noise model can only \emph{add} probabilities; it can never
make them cancel, so it has no slot for interference and is silent about the
pre-measurement state. We say this plainly in \cref{sec:channel}, and we keep
the scope honest in \cref{sec:scope}: the contribution is forward budgeting
under per-bit independence, complementary to (not competitive with) inverse
mitigation, and blind by construction to measurement crosstalk.

\Cref{sec:channel} sets up the Bernoulli readout channel. \Cref{sec:calculus}
derives the forward $(\eps,\omega)$ calculus. \Cref{sec:worked} works the
syndrome-bit example and tabulates the linearization error. \Cref{sec:related}
positions the work. \Cref{sec:scope} states the limitations, and
\cref{sec:conclusion} concludes.


\section{The Bernoulli readout channel}
\label{sec:channel}

\subsection{The atom: a measured bit as an asymmetric channel}

Let a qubit carry a definite latent classical value $t\in\bits$ in the
measurement (computational) basis, and let $b\in\bits$ be the observed readout.
We model the measurement as a memoryless binary channel.

\begin{definition}[Bernoulli readout channel]
\label{def:atom}
A \emph{Bernoulli readout channel} with parameters $(\eps,\omega)\in[0,1]^2$ is
the $2\times2$ row-stochastic matrix
\begin{equation}
\label{eq:channel}
  Q \;=\;
  \begin{pmatrix} \Prob(b=0\given t=0) & \Prob(b=1\given t=0) \\[2pt]
                  \Prob(b=0\given t=1) & \Prob(b=1\given t=1) \end{pmatrix}
  \;=\;
  \begin{pmatrix} 1-\eps & \eps \\[2pt] \omega & 1-\omega \end{pmatrix},
\end{equation}
where $\eps = \Prob(b=1\given t=0)$ is the \emph{false-positive rate} and
$\omega = \Prob(b=0\given t=1)$ is the \emph{false-negative rate}.
\end{definition}

Rows sum to one, so $Q$ is a valid classical channel for every
$(\eps,\omega)\in[0,1]^2$. Two special cases recur.

\begin{itemize}[leftmargin=1.4em]
  \item \textbf{Symmetric channel} ($\eps=\omega=q$): the binary symmetric
        channel, the implicit model behind single-rate readout error.
  \item \textbf{Z-channel} ($\omega=0$, $\eps>0$): only $0\!\to\!1$ flips occur;
        a read $0$ is certainly correct. Its dual ($\eps=0$, $\omega>0$) admits
        only $1\!\to\!0$ flips, so a read $1$ is certainly correct.
\end{itemize}

\begin{remark}[The physics fixes the sign of the asymmetry]
\label{rem:t1}
For superconducting qubits the measurement window is long enough that $T_1$
relaxation of the excited state contributes to errors, biasing misreads toward
$\ket{1}\!\to\!\ket{0}$ \citep{tannu2019mitigating,hicks2021readout}. As a
concrete instance, \citet{tannu2019mitigating} report an IBM device reading the
all-zero state at $84\%$ fidelity but the all-one state at only $62\%$, precisely
the relaxation-driven $\ket{1}\!\to\!\ket{0}$ bias. In our
convention that is the false-negative direction, so readout typically sits in the
regime $\omega > \eps$, partway toward the dual Z-channel. The calculus below
never assumes $\eps=\omega$; preserving the inequality is the point.
\end{remark}

\subsection{Bayesian recovery and why the latent value is well posed}

Given a prior $\pi=\Prob(t=1)$ and an observation $b$, Bayes' rule recovers the
posterior over the latent value, e.g.
\begin{equation}
\label{eq:bayes}
  \Prob(t=1\given b=1) \;=\;
  \frac{(1-\omega)\,\pi}{(1-\omega)\,\pi + \eps\,(1-\pi)},
  \qquad
  \Prob(t=1\given b=0) \;=\;
  \frac{\omega\,\pi}{\omega\,\pi + (1-\eps)(1-\pi)} .
\end{equation}
This is the inferential (inverse) reading of the atom; classification-quality
summaries such as positive/negative predictive value follow from it and are the
subject of separate framework work \citep{bernoulliMeasures}. Our concern is the
\emph{forward} direction: pushing $(\eps,\omega)$ through computation.

\paragraph{Faithfulness is a hypothesis, stated honestly.}
\Cref{def:atom} presumes a single joint distribution over the latent value $t$
and the readout $b$. This is faithful precisely when the qubit's density matrix
$\rho$ is diagonal in the measurement basis, an eigenstate, a classical
mixture, or a fully decohered state, which is exactly the readout sector. When
$\rho$ has off-diagonal coherences the model is silent: a latent-value-plus-noise
description can only add probabilities and can never reproduce interference. We
therefore make no claim about pre-measurement amplitudes, amplitude
amplification, or entangled-state correlations; the channel of \cref{def:atom}
is a statement about the bits a device \emph{reports}, not about the quantum
state it holds. Within that diagonal sector the model is exact, and the calculus
that follows is exact with it.

\subsection{Multi-qubit readout and the assignment matrix}

For $n$ qubits read \emph{independently}, with qubit $i$ carrying channel
$Q_i=\bigl(\begin{smallmatrix}1-\eps_i & \eps_i\\ \omega_i & 1-\omega_i\end{smallmatrix}\bigr)$,
the joint channel over the $2^n$ truth assignments and $2^n$ readout patterns is
the Kronecker product
\begin{equation}
\label{eq:kron}
  A \;=\; Q_1 \tensor Q_2 \tensor \cdots \tensor Q_n .
\end{equation}
Equation~\eqref{eq:kron} is the \emph{assignment matrix} $A$ of readout-error
mitigation, in its tensored/local form. We stress at the outset (and return to
it in \cref{sec:related}) that the factorization~\eqref{eq:kron} is known
practice; our contribution is not the factorization but the symbolic forward
\emph{calculus} we now build on top of it, which never needs to materialize or
invert the $2^n\times2^n$ matrix.


\section{The forward \texorpdfstring{$(\eps,\omega)$}{(epsilon, omega)} calculus}
\label{sec:calculus}

Throughout this section, ``bit $b_i$'' denotes a measured bit drawn through an
independent Bernoulli readout channel with rates $(\eps_i,\omega_i)$ over a
latent value $t_i$, and we compute the rates of a derived bit
$g(b_1,\dots,b_n)$ relative to its derived latent value $g(t_1,\dots,t_n)$. We
assume the channels are mutually independent, and that each measured bit enters
$g$ \emph{once}, so the derived bit is a fan-out-free Boolean \emph{tree} of
distinct inputs. Reusing one readout in several places (reconvergent fan-out)
correlates those inputs and voids the product and parity rules below: $b\wedge b
= b$ has false-positive rate $\eps$, not the $\eps^2$ the independent rule would
report. We flag this structural caveat alongside the physical one in
\cref{sec:scope}. The reader should read $\eps$ as
``error when the true answer is $0$ (a false alarm)'' and $\omega$ as ``error
when the true answer is $1$ (a miss)'' uniformly across all derived bits. A
subscript on a rate is either a qubit index ($\eps_i$) or a connective tag
($\eps_\wedge$, $\eps_\vee$, $\omega_\wedge$); context fixes which.

It is convenient to also track the \emph{conditional flip probabilities} given
the latent value, which for a single bit are exactly $\eps$ (when $t=0$) and
$\omega$ (when $t=1$).

\subsection{Negation}

\begin{proposition}[\textsc{not}]
\label{prop:not}
If $b$ has rates $(\eps,\omega)$ for latent value $t$, then $\lnot b$ has rates
$(\eps',\omega')=(\omega,\eps)$ for latent value $\lnot t$.
\end{proposition}

\begin{proof}
The latent value of $\lnot b$ is $\lnot t$. A false positive for $\lnot b$ is the
event $\lnot b=1$ while $\lnot t=0$, i.e. $b=0$ while $t=1$, which occurs with
probability $\omega$. Symmetrically the false-negative rate of $\lnot b$ is
$\eps$. Negation \emph{swaps} the two rates; on the symmetric channel it would be
invisible, which is exactly why a single-rate model cannot see it.
\end{proof}

\Cref{prop:not} also pins down the duality used below: \textsc{or} is
\textsc{not}-\textsc{and}-\textsc{not}, so every \textsc{and} identity has an
\textsc{or} dual obtained by swapping $\eps\leftrightarrow\omega$ and the roles
of $0$ and $1$.

\subsection{Conjunction and disjunction}

\begin{theorem}[\textsc{and} and \textsc{or}]
\label{thm:andor}
Let $b_1,\dots,b_n$ be independent measured bits with rates
$(\eps_i,\omega_i)$ for latent values $t_i$.
\begin{enumerate}[label=(\roman*),leftmargin=2em]
  \item The conjunction $c=\bigwedge_i b_i$ has latent value $\bigwedge_i t_i$.
        Conditioned on all latent inputs being false ($t_i=0$ for every $i$), its
        false-positive rate is
        \begin{equation}
        \label{eq:and-fpr}
          \eps_{\wedge} \;=\; \prod_{i=1}^{n} \eps_i ;
        \end{equation}
        this product is the all-negative corner of the general
        latent-assignment-weighted rate, the smallest such corner rather than a
        worst case (see the proof). Its false-negative rate, which conditions on
        the only latent assignment making the conjunction true (all $t_i=1$), is
        \begin{equation}
        \label{eq:and-fnr}
          \omega_{\wedge} \;=\; 1 - \prod_{i=1}^{n}(1-\omega_i).
        \end{equation}
  \item Dually, the disjunction $d=\bigvee_i b_i$ has latent value
        $\bigvee_i t_i$, with
        \begin{equation}
        \label{eq:or}
          \omega_{\vee} \;=\; \prod_{i=1}^{n}\omega_i,
          \qquad
          \eps_{\vee} \;=\; 1 - \prod_{i=1}^{n}(1-\eps_i).
        \end{equation}
\end{enumerate}
\end{theorem}

\begin{proof}
\emph{(i), false positives.} A false positive for the conjunction is
$c=1$ while the latent conjunction $\bigwedge_i t_i=0$. Condition on the latent
assignment. The event $c=1$ requires every $b_i=1$; if some $t_j=0$ then
producing $b_j=1$ costs a factor $\eps_j$, while producing $b_i=1$ for the other
indices costs $\eps_i$ when $t_i=0$ and $1-\omega_i$ when $t_i=1$. When
\emph{all} $t_i=0$ this product is $\prod_i \eps_i$, giving~\eqref{eq:and-fpr}.
That all-negative corner is the \emph{smallest}, not the largest, such rate: an
input whose latent value is true replaces its $\eps_i$ by $1-\omega_i$ and so
raises the conjunction's false-positive rate toward the rates of the remaining
inputs. The unconditional rate is the cardinality-weighted average of these
corners derived for the intersection of approximate sets by
\citet{bernoulliComposition}; we quote the all-negative corner because it is the
case that governs a conjunction of independently-false readouts. \emph{(i), false negatives.} A false
negative is $c=0$ while $\bigwedge_i t_i=1$, i.e. all $t_i=1$ but at least one
$b_i=0$. With all $t_i=1$ the bits flip independently to $0$ with probability
$\omega_i$, so $\Prob(c=0)=1-\prod_i(1-\omega_i)$, giving~\eqref{eq:and-fnr}.
\emph{(ii)} follows from (i) by \cref{prop:not} and De~Morgan, swapping
$\eps\leftrightarrow\omega$.
\end{proof}

\begin{remark}[Conjunction is an intersection of approximate sets]
\Cref{eq:and-fpr} is the multiplicative false-positive rule for the intersection
of independent random approximate sets in Bernoulli type theory: the membership
predicate of $A\cap B$ is the \textsc{and} of the two membership predicates, and
on elements absent from both sets its false-positive rate is the product
$\eps_A\eps_B$ \citep{bernoulliSets,bernoulliComposition}. Over the full query
population the rate is the three-region weighted average of
\citet{bernoulliComposition}, of which $\eps_A\eps_B$ is the both-absent corner;
\cref{thm:andor} quotes that corner. A Bloom-filter conjunction is the
$\omega_i=0$ specialization; readout-derived conjunctions live at the asymmetric
interior. The calculus is one object serving both.
\end{remark}

\subsection{Chains: the composition theorem}

A particularly clean corollary governs any one-sided ``chain'' of independent
error events, a value that is wrong if \emph{any} stage errs. Let stage $i$
have error probability $\eta_i$ (a single rate, on the relevant side).

\begin{corollary}[Composition theorem]
\label{cor:composition}
For independent stages with per-stage error probabilities $\eta_1,\dots,\eta_n$,
the probability that the chain errs is
\begin{equation}
\label{eq:composition}
  \eta_{\mathrm{total}} \;=\; 1 - \prod_{i=1}^{n}\bigl(1-\eta_i\bigr).
\end{equation}
\end{corollary}

\begin{proof}
The chain is correct iff every stage is correct; by independence this has
probability $\prod_i(1-\eta_i)$, and~\eqref{eq:composition} is its complement.
\end{proof}

Equation~\eqref{eq:composition} is the false-negative rule for conjunction
\eqref{eq:and-fnr} (and, dually, the false-positive rule for disjunction): a
conjunction is a false negative iff \emph{any} conjunct is independently missed.
It is also exactly the composition rule for chained Bernoulli maps and cipher
maps, where it governs the total error of a composed approximate computation. For
small rates it linearizes to $\eta_{\mathrm{total}}\approx\sum_i\eta_i$, the
familiar union-bound estimate; we will see in \cref{sec:worked} that the
linearization is one-sided and can mislead.

\subsection{Parity}

The most important derived bit for quantum error correction is a \emph{parity}.
Here the exact answer is cleanest and the gap from common practice is largest.

\begin{theorem}[Parity]
\label{thm:parity}
Let $X_1,\dots,X_n$ be independent $\bits$-valued bits with
$\Prob(X_i=1)=q_i$. Then
\begin{equation}
\label{eq:parity}
  \Prob\!\Bigl(\bigxor_{i=1}^{n} X_i = 1\Bigr)
  \;=\; \frac12\Bigl(1 - \prod_{i=1}^{n}(1-2q_i)\Bigr).
\end{equation}
\end{theorem}

\begin{proof}
Encode each bit by the $\{\pm1\}$ variable $S_i=1-2X_i$, so $S_i=+1$ when
$X_i=0$ and $S_i=-1$ when $X_i=1$, with $\mathbb{E}[S_i]=1-2q_i$. Parity in
$\bits$ becomes a product in $\{\pm1\}$: $S=\prod_i S_i = 1-2\bxor_i X_i$. By
independence $\mathbb{E}[S]=\prod_i\mathbb{E}[S_i]=\prod_i(1-2q_i)$. Since
$\Prob(\bxor_i X_i=1)=\Prob(S=-1)=\tfrac12(1-\mathbb{E}[S])$,
equation~\eqref{eq:parity} follows.
\end{proof}

Theorem~\ref{thm:parity} is classical: the $\pm1$/Fourier parity-bias identity is
standard in the analysis of Boolean functions \citep{odonnell2014}, and it
underlies the noisy-Boolean-formula reliability tradition of \citet{vonneumann1956}
and \citet{pippenger1988}. We now state its asymmetric,
\emph{conditioned} form, which is what a syndrome bit actually needs and what a
symmetric single-rate treatment cannot express.

\begin{corollary}[Asymmetric parity of measured bits]
\label{cor:parity-asym}
Let $b_1,\dots,b_n$ be independent measured bits with rates $(\eps_i,\omega_i)$
over latent values $t_i$, and let the derived bit be the parity
$p=\bxor_i b_i$ with latent value $\tau=\bxor_i t_i$. Conditioned on the latent
assignment $(t_1,\dots,t_n)$, the per-bit flip probability is $r_i=\eps_i$ if
$t_i=0$ and $r_i=\omega_i$ if $t_i=1$, and the probability that the measured
parity disagrees with $\tau$ is
\begin{equation}
\label{eq:parity-asym}
  \Prob\bigl(p \neq \tau \given t_1,\dots,t_n\bigr)
  \;=\; \frac12\Bigl(1 - \prod_{i=1}^{n}\bigl(1-2r_i\bigr)\Bigr).
\end{equation}
In particular the parity flips iff an \emph{odd} number of its bits flip, so the
two ``sides'' (true parity $0$ vs.\ true parity $1$) generally carry different
error, and~\eqref{eq:parity-asym} is the rate that must be reported for each
latent assignment.
\end{corollary}

\begin{proof}
Apply \cref{thm:parity} to the \emph{flip} indicators $F_i=\Ind{b_i\neq t_i}$,
which are independent with $\Prob(F_i=1)=r_i$ as conditioned. The parity disagrees
with $\tau$ iff $\bxor_i F_i=1$, and~\eqref{eq:parity-asym} is
equation~\eqref{eq:parity} with $q_i\mapsto r_i$.
\end{proof}

\begin{remark}[The single-rate special case]
When all $r_i=q$, equation~\eqref{eq:parity-asym} collapses to
$\tfrac12(1-(1-2q)^n)$, the symmetric parity error. The asymmetric form keeps the
distinction between $\eps$-flips and $\omega$-flips alive through the parity tree;
this is precisely the bookkeeping the classical single-rate identities, and the
symmetrized device tooling, throw away.
\end{remark}

\subsection{Kronecker / tensor composition over independent qubits}

Finally, the multi-qubit story. The forward calculus composes over independent
qubits exactly as the assignment matrix~\eqref{eq:kron} factorizes.

\begin{proposition}[Tensor factorization of the forward calculus]
\label{prop:tensor}
For independent qubits the joint readout channel is $A=\bigotimes_i Q_i$, and
the rate of any Boolean function $g$ of the readouts depends on the $Q_i$ only
through the closed forms of \cref{prop:not,thm:andor,thm:parity}; no entry of
the $2^n\times2^n$ matrix $A$ need be formed. Equivalently, the marginal channel
on any subset $S$ of qubits is $\bigotimes_{i\in S} Q_i$, and the rate of a
derived bit reading only those qubits is computed from $\{(\eps_i,\omega_i)\}_{i\in S}$
alone.
\end{proposition}

\begin{proof}
Independence makes the joint distribution of the readouts the product of the
marginals, which is the statement $A=\bigotimes_i Q_i$. Every derived rate in
this section was computed from marginal flip probabilities under independence, so
it is a function of the $(\eps_i,\omega_i)$ and not of off-diagonal joint
structure (of which there is none under independence).
\end{proof}

\Cref{prop:tensor} is the typed, symbolic shadow of tensored mitigation: the same
factorization the mitigation literature exploits numerically, read forward and in
closed form. The next section shows what it buys.

\subsection{An operator view: assignment matrices and the readout spectrum}
\label{sec:operator}

It is worth recasting the calculus in the operator language the
quantum-computing audience already uses, because two of its central objects are
this paper's objects under another name. The recasting is bookkeeping, not new
theorems; its value is that it makes the parity result visibly the meeting point
of two existing traditions.

\paragraph{The dictionary.} The single-bit readout channel of
\cref{def:atom} carries exactly the information in the \emph{assignment matrix}
(confusion matrix) that readout-error mitigation calibrates and inverts
\citep{nation2021m3,bravyi2021mitigating}, differing from it only by the
transpose convention. Writing it in the column-stochastic convention the
mitigation literature prefers, with the observed distribution as a column vector
acted on from the left, it is
\begin{equation}
\label{eq:assignment}
  A_i \;=\;
  \begin{pmatrix} 1-\eps_i & \omega_i \\[2pt] \eps_i & 1-\omega_i \end{pmatrix},
\end{equation}
the transpose of \cref{eq:channel}. For $n$ qubits read independently the joint
assignment matrix is the Kronecker product $\bigotimes_i A_i$, the transpose of
the row-stochastic $A$ of \cref{eq:kron}: the same tensored (local) assignment
matrix used in practice, written in the column convention. The element-wise independence that makes $A$ factor is the same
independence that lets the forward calculus of \cref{sec:calculus} factor through
the per-qubit rates. So the field's assignment matrix is this paper's channel,
and ``tensored mitigation'' is the Kronecker factorization the forward calculus
already exploits; the only difference is direction. Mitigation inverts $A$ to
pull an ideal distribution back from observed counts; the present calculus pushes
$A$ \emph{forward} through Boolean post-processing, never forming or inverting the
$2^n\times2^n$ matrix.

\paragraph{The spectrum.} The parity identity has a clean spectral reading in
this language. For the symmetric per-bit flip channel ($\eps_i=\omega_i=q_i$) the
$2\times2$ matrix \eqref{eq:assignment} has eigenvalues $1$ and $1-2q_i$: the
eigenvalue $1$ belongs to the uniform (stationary) mode $(\tfrac12,\tfrac12)$,
and the nontrivial eigenvalue $1-2q_i$ belongs to the difference, or
\emph{parity}, mode $(\tfrac12,-\tfrac12)$. That single number $1-2q_i$ is three
things at once: (i) the channel's nontrivial eigenvalue; (ii) the $\pm1$ bias
$\mathbb{E}[(-1)^{F_i}]=\mathbb{E}[S_i]$ that the proof of \cref{thm:parity}
already rides on, with $S_i=1-2F_i$; and (iii) the Walsh--Hadamard (Fourier)
coefficient of the bit's flip distribution on the parity character. Under the
Kronecker product these per-bit eigenvalues multiply on the all-bits parity mode,
so the syndrome/parity identity
\begin{equation}
\label{eq:parity-spectral}
  \Prob\bigl(\text{parity wrong}\bigr)
  \;=\; \frac12\Bigl(1 - \prod_{i=1}^{n}(1-2q_i)\Bigr)
\end{equation}
is exactly one half of one minus the product of the channels' nontrivial
eigenvalues: the spectral coefficient of the tensored readout channel on the
parity mode. This is the same $\pm1$-bias product the classical
noisy-Boolean-formula tradition computes \citep{vonneumann1956,pippenger1988},
the order-one Walsh coefficient in the modern Fourier framing \citep{odonnell2014}.
Read this way, \cref{thm:parity} is visibly where the quantum operator spectrum
of the assignment matrix and the classical Fourier analysis of noisy Boolean
formulas coincide; the asymmetric \cref{cor:parity-asym} is the same coefficient
with the symmetric bias $1-2q_i$ replaced by the latent-conditioned flip bias
$1-2r_i$, $r_i\in\{\eps_i,\omega_i\}$. Only in the symmetric case does this bias
coincide with the channel's nontrivial eigenvalue; the asymmetric assignment
matrix~\eqref{eq:assignment} has eigenvalues $1$ and $1-\eps_i-\omega_i$, so it
is the $\pm1$ bias, not that eigenvalue, that the parity formula carries.

One honesty clause keeps the framing from being mistaken for a coherence claim.
The bra-ket and operator language here is ordinary finite-dimensional linear
algebra: the vectors are classical probability distributions, not amplitudes,
there is no Born-rule square, and the entire discussion lives on the diagonal of
the density matrix, the readout sector of \cref{sec:channel}. The eigenvalue
$1-2q_i$ is a relaxation rate of a stochastic matrix, not a quantum phase.


\section{Worked example: the reliability of a syndrome bit}
\label{sec:worked}

\subsection{Setup}

In a stabilizer quantum error-correcting code, each \emph{syndrome bit} is the
parity of a set of ancilla measurements; a nonzero syndrome flags a detected
error, and decoders consume the syndrome to infer a correction
\citep{fowler2012surface}. The reliability of a syndrome bit is therefore the
reliability of a parity of noisy measured bits, precisely
\cref{cor:parity-asym}. Take $n$ ancillas, each read through an independent
channel, and consider the symmetric per-measurement rate $q_i=q$ for clarity, so
that $r_i=q$ for every $i$ in \cref{cor:parity-asym} (which handles the asymmetric
and inhomogeneous case verbatim).
The exact syndrome-flip probability is
\begin{equation}
\label{eq:syndrome-exact}
  P_{\mathrm{flip}}^{\mathrm{exact}}(n,q)
  \;=\; \frac12\Bigl(1-(1-2q)^n\Bigr).
\end{equation}

\subsection{The linearization the field uses, and its error}

A common back-of-envelope estimate treats the syndrome error as additive in the
number of measurements,
\begin{equation}
\label{eq:syndrome-lin}
  P_{\mathrm{flip}}^{\mathrm{lin}}(n,q) \;=\; n\,q,
\end{equation}
the union-bound / ``about $n$ times the per-measurement error'' rule. This is the
first-order term of~\eqref{eq:syndrome-exact}: expanding $(1-2q)^n$ gives
$P^{\mathrm{exact}} = nq - \binom{n}{2}(2q)^2/2 + \cdots = nq - n(n-1)q^2 + O(q^3)$,
so the linear rule \emph{overestimates} the true error, and ignores the
saturation of~\eqref{eq:syndrome-exact} toward $\tfrac12$ as $nq$ grows
(a parity of many fair-ish bits is itself nearly fair, never worse than a coin).
For $nq \gtrsim \tfrac12$ the linear rule can exceed $1$ and is meaningless,
whereas the exact value is bounded by $\tfrac12$.

\Cref{tab:syndrome} tabulates both, with the relative overestimate
$\bigl(P^{\mathrm{lin}}-P^{\mathrm{exact}}\bigr)/P^{\mathrm{exact}}$.

\begin{table}[t]
\centering
\caption{Exact syndrome-bit flip probability~\eqref{eq:syndrome-exact} versus the
linear ``$n$ times per-measurement error'' rule~\eqref{eq:syndrome-lin}, for a
parity of $n$ ancillas at per-measurement error $q$. The linear rule always
overestimates; the overestimate grows with $n$ and $q$, and the linear value
becomes unphysical ($>0.5$, even $\ge1$) while the exact value saturates at
$0.5$.}
\label{tab:syndrome}
\small
\begin{tabular}{c c r r r}
\toprule
$q$ & $n$ & $P^{\mathrm{exact}}$ & $P^{\mathrm{lin}}=nq$ & rel.\ overestimate \\
\midrule
0.01 & 2 & 0.019800 & 0.020000 & \phantom{0}1.0\% \\
0.01 & 4 & 0.038816 & 0.040000 & \phantom{0}3.1\% \\
0.01 & 8 & 0.074618 & 0.080000 & \phantom{0}7.2\% \\
\midrule
0.05 & 2 & 0.095000 & 0.100000 & \phantom{0}5.3\% \\
0.05 & 4 & 0.171950 & 0.200000 & 16.3\% \\
0.05 & 8 & 0.284766 & 0.400000 & 40.5\% \\
\bottomrule
\end{tabular}
\end{table}

\subsection{Reading the table}

At the low end ($q=0.01$, $n=2$) the linear rule is off by one percent, which is
why it survives as folklore. But the error compounds: at $q=0.05$ over an
$8$-ancilla parity the linear rule reports $0.40$ against a true $0.285$, a
$40.5\%$ relative overestimate, and at slightly larger $nq$ it crosses $0.5$ and
becomes nonsensical while the truth never can. The exact form costs nothing, it
is one product, and it is the form that the asymmetric \cref{cor:parity-asym}
generalizes to inhomogeneous and $T_1$-biased ancillas (different $\eps_i,\omega_i$
per qubit, e.g. when some ancillas are prepared/measured in the $\ket{1}$-heavy
regime). The point is not that practitioners cannot compute~\eqref{eq:syndrome-exact};
it is that there is a reusable \emph{calculus} that delivers it, and the
matching \textsc{and}/\textsc{or}/chain forms for majority and feed-forward
post-processing, directly from the per-qubit rates, asymmetry intact, without
materializing the assignment matrix.

\subsection{The asymmetric case: a number the symmetric rule cannot give}
\label{sec:worked-asym}

\Cref{tab:syndrome} is symmetric ($q_i=q$), so every value in it is reproducible
from the single-rate identity~\eqref{eq:parity} we attribute to prior art; the
\emph{asymmetric} content of \cref{cor:parity-asym} never appears. We now make it
appear. Take a four-ancilla syndrome with $T_1$-biased, inhomogeneous channels:
qubits $1,2$ at $(\eps_i,\omega_i)=(0.02,0.10)$ and qubits $3,4$ at $(0.04,0.12)$,
each with $\omega_i>\eps_i$ as relaxation dictates. By \cref{cor:parity-asym} the
syndrome-flip probability depends on the \emph{latent assignment}
$(t_1,\dots,t_4)$ through the per-bit flip rate $r_i=\eps_i$ when $t_i=0$ and
$r_i=\omega_i$ when $t_i=1$, not merely on the true syndrome $\tau=\bxor_i t_i$.
\Cref{tab:syndrome-asym} evaluates it on four assignments and contrasts them with
the symmetrized single-rate prediction~\eqref{eq:syndrome-exact} at the mean rate
$q=\bar r=0.07$, the best a one-rate model can do.

\begin{table}[t]
\centering
\caption{Asymmetric syndrome-flip probability~\eqref{eq:parity-asym} for a parity
of four $T_1$-biased ancillas, qubits $1,2$ at $(\eps,\omega)=(0.02,0.10)$ and
qubits $3,4$ at $(0.04,0.12)$. The error tracks the latent assignment, not just the
true syndrome $\tau$: the all-excited word $1111$ is misread almost three times as
often as the all-ground word $0000$ though both carry $\tau=0$. A symmetrized
single-rate model reports one value ($0.226$) for every assignment, overstating the
cold corner roughly twofold and understating the hot one.}
\label{tab:syndrome-asym}
\small
\begin{tabular}{c c r}
\toprule
latent word $t_1t_2t_3t_4$ & true syndrome $\tau$ & exact flip prob.\ \eqref{eq:parity-asym} \\
\midrule
$0000$ & $0$ & 0.110 \\
$1000$ & $1$ & 0.175 \\
$0011$ & $0$ & 0.234 \\
$1111$ & $0$ & 0.315 \\
\midrule
symmetrized single rate ($q=0.07$) & any & 0.226 \\
\bottomrule
\end{tabular}
\end{table}

The three $\tau=0$ rows ($0.110$, $0.234$, $0.315$) already span a factor of nearly
three, ordered by how many ancillas sit in the relaxation-prone $\ket{1}$ state,
and the single-rate column collapses all of it to $0.226$, wrong on \emph{both}
sides at once. This is the bookkeeping \cref{cor:parity-asym} keeps and a symmetric
treatment throws away, computed from the eight per-qubit rates alone with no
$2^4\times2^4$ assignment matrix in sight.

\begin{remark}[Majority vote and feed-forward, in one line each]
A \emph{majority} of $2m+1$ independent measured copies of one latent bit has
error $\sum_{k=m+1}^{2m+1}\binom{2m+1}{k} r^k (1-r)^{2m+1-k}$ on the conditioned
side with flip rate $r\in\{\eps,\omega\}$, the Bernoulli amplification rule,
which (unlike coherent amplitude amplification) requires a stable resampleable
latent value, available precisely because we are in the diagonal readout sector.
A \emph{feed-forward conditional} ``apply $G$ if $b=1$'' fires erroneously with
probability $\eps$ when it should not and is skipped erroneously with probability
$\omega$ when it should fire; these compose by \cref{cor:composition} down a
chain of conditionals. Both are immediate instances of the calculus.
\end{remark}


\section{Related work and positioning}
\label{sec:related}

We situate the calculus against five neighbors, citing each and stating the
distinction. We claim neither the single-bit asymmetric channel nor the
independent-qubit tensor factorization as original; the claim is the reusable
\emph{forward asymmetric Boolean calculus}, its readout application, and the
unification with classical probabilistic data structures.

\paragraph{(a) Matrix-inversion readout mitigation (the inverse paradigm).}
M3 \citep{nation2021m3}, the multiqubit treatment of
\citet{bravyi2021mitigating}, and tensored/local assignment-matrix mitigators
build the assignment matrix $A$ (our~\eqref{eq:kron}) and approximately invert it
to recover an ideal distribution from observed counts. This is an \emph{inverse},
\emph{numerical}, \emph{per-circuit} procedure operating on measured data. The
work nearest our motivating examples is the mid-circuit-measurement and
feed-forward readout mitigation of \citet{koh2024readout}, which targets exactly
the derived and fed-forward readouts we post-process; it too is an inverse
mitigation of measured outcomes, not a forward symbolic rate calculus. Our
calculus is its forward complement: given the rates, predict and budget the error
on a specific derived bit, symbolically, before data. The tensored assignment
matrix is the linear-algebra shadow of our \cref{prop:tensor}; we add the
symbolic forward layer over it and never invert anything. The two are
complementary, not competing, one designs the budget, the other cleans the
data.

\paragraph{(b) Effectus theory / Adams--Jacobs type theory.}
Effectus theory and the Adams--Jacobs type theory for probabilistic and Bayesian
reasoning \citep{cho2015effectus,adams2015typetheory} give a categorical
``algebra of predicates'' in which predicates form effect algebras and Boolean
and probabilistic logic appear as special cases of the quantum case, the same
structural DNA as a typed calculus of predicates under noise. It is, however, a
\emph{foundation for the structure} of probabilistic and quantum theory and for
Bayesian conditioning; it is not a device readout-error model and carries no
$(\eps,\omega)$ rate-propagation rules. Our contribution is the concrete
asymmetric rate calculus and its quantum-readout instantiation, which effectus
theory neither provides nor subsumes. This is the most important neighbor to
distinguish, and the distinction is precisely \emph{rates}: we propagate numbers,
not just predicate structure.

\paragraph{(c) Classical parity / noisy-Boolean-formula identities.}
The exact $\pm1$/Fourier parity-bias identity~\eqref{eq:parity} is standard in the
analysis of Boolean functions \citep{odonnell2014}, while the reliability of noisy
Boolean formulas under a single symmetric flip rate is the tradition of
\citet{vonneumann1956} and \citet{pippenger1988}; the analysis of
XOR-PUFs uses the same product form. That literature owns the parity formula, but
uses a \emph{single symmetric} flip rate. Our originality on this axis is the
asymmetric two-rate bookkeeping (\cref{prop:not,thm:andor,cor:parity-asym}), the
quantum-measurement application, and the unification, not the parity formula
itself, which we cite as prior art.

\paragraph{(d) Asymmetric quantum hypothesis testing / channel certification.}
A two-sided false-positive/false-negative model of a \emph{single} measurement
already appears in asymmetric quantum hypothesis testing and channel
certification, for instance the false-negative-free certification of
\citet{krawiec2021certification}. This is the nearest neighbor of our atom
(\cref{def:atom}) and we cite it to preempt ``this is just type~I/II error.'' The
difference is the \emph{composition layer}: those works study one decision, while
our contribution is the closed-form propagation of the two rates through Boolean
combinations and tensor composition of \emph{many} measured bits.

\paragraph{(e) Asymmetry-from-relaxation physics.}
\citet{tannu2019mitigating} establish and exploit state-dependent (asymmetric)
readout bias from $T_1$ relaxation, with measured errors of roughly $8$--$30\%$;
\citet{hicks2021readout} use the same asymmetry via readout rebalancing. These
motivate \cref{def:atom} and \cref{rem:t1}, they are the source of the $\omega>\eps$
regime our calculus is built to preserve. They mitigate; we propagate.

\paragraph{Unification with classical probabilistic data structures.}
Finally, the calculus is not new \emph{to us}: it is the error model of Bernoulli
type theory, the framework of random approximate sets, maps, and relations whose
special cases include Bloom filters \citep{bloom1970}, perfect hash filters, and
cipher maps \citep{bernoulliSets,bernoulliComposition,bernoulliMaps,bernoulliRelations}.
For self-containedness the unification needs no external reference: a random
approximate set is the membership indicator of a set $A$, reported through an
independent per-element asymmetric channel~\eqref{eq:channel} with false-positive
rate $\eps$ on nonmembers and false-negative rate $\omega$ on members, and set
operations compose those two rates by the closed forms of \cref{sec:calculus}. The
companion papers develop the consequences, but the unification claim rests only on
this one-line model.
A Bloom filter is the one-sided ($\omega=0$) corner; a $T_1$-biased syndrome
ancilla is an asymmetric interior point; both obey the same
\cref{prop:not,thm:andor,thm:parity,cor:composition}. Recognizing quantum readout
as an instance of this model is what lets a single calculus serve a classical
data structure and a quantum-hardware error budget at once.


\section{Scope and limitations}
\label{sec:scope}

We state the boundaries plainly, to forestall overclaiming.

\begin{enumerate}[leftmargin=1.6em]
  \item \textbf{Forward, not inverse.} The contribution predicts and budgets the
        error of a derived bit from given per-qubit rates. It does \emph{not}
        recover an ideal distribution from observed counts and is \emph{not} a
        competitor to M3 or other mitigation; it is their design-time complement
        \citep{nation2021m3,bravyi2021mitigating}.

  \item \textbf{Independence is assumed, and is the reference point.} Every closed
        form above holds under per-bit channel independence, equivalently, the
        tensored assignment matrix~\eqref{eq:kron}. Real devices exhibit
        \emph{measurement crosstalk}: the readout of one qubit is correlated with
        the state or readout of its neighbors, which is why correlated models such
        as the continuous-time Markov process (CTMP) model of
        \citet{bravyi2021mitigating} exist. Independence can also fail
        \emph{structurally} rather than physically: a single readout reused across
        several inputs of the derived bit (reconvergent fan-out, \cref{sec:calculus})
        correlates those inputs exactly as crosstalk does, so a derived bit that is
        not a fan-out-free tree needs the same non-factorized treatment. The
        independent prediction is best read
        as the correlation-free \emph{reference}: it is exact given the per-qubit
        marginals alone, and a correlated model perturbs it in a sign-determined
        way rather than by a uniform bias. The sign is legible on a parity, which
        flips only under an odd number of bit flips: positively correlated misreads
        make coincident flips more likely and so push the parity error \emph{below}
        the independent value, while negatively correlated misreads push it
        \emph{above}. The independent rate is thus a two-sided anchor, not a worst
        case, and bounding the perturbation for a given CTMP generator is the
        natural next step. A correlated-channel extension, propagating a joint,
        non-factorized readout model through the same Boolean operations, is future
        work; the present calculus is the independent/tensor regime.

  \item \textbf{Diagonal (readout) sector only.} The latent-value model is exact
        precisely where the density matrix is diagonal in the measurement basis.
        It says nothing about pre-measurement coherences, interference, amplitude
        amplification, or entangled-state correlations; for those the
        single-joint-distribution premise fails (this is the content of Bell /
        Kochen--Specker / Fine obstructions, outside our scope here). Our claims
        concern the bits a device reports, full stop.

  \item \textbf{Originality, precisely.} Not claimed: the single-bit asymmetric
        channel (\cref{def:atom}); the tensor factorization of independent-qubit
        readout (\cref{eq:kron}); the parity formula~\eqref{eq:parity}. Claimed:
        the reusable asymmetric forward Boolean calculus carrying \emph{both}
        rates (\cref{sec:calculus}); its application to quantum readout and the
        classical post-processing of measurement outcomes; and the unification of
        a quantum-hardware error model with the Bernoulli error model of classical
        probabilistic data structures.

  \item \textbf{Empirical validation.} This is a methods/perspective paper; the
        closed forms are derived and the syndrome tables are computed exactly from
        them. A hardware or simulator study comparing predicted derived-bit rates
        against measured ones, and quantifying the independence-violation gap from
        crosstalk, would strengthen the claims but is out of scope for this draft.
        We make no empirical claim beyond the exactness of the derivations under
        the stated independence assumption.
\end{enumerate}


\section{Conclusion}
\label{sec:conclusion}

Quantum readout error is a classical asymmetric bit-flip channel on the diagonal
sector of the density matrix, and the bits a quantum program reports are Boolean
functions of such channels. We have given the closed-form, asymmetric, forward
calculus that propagates a false-positive rate $\eps$ and a false-negative rate
$\omega$ through negation, conjunction, disjunction, parity, chained composition,
and independent-qubit tensoring, carrying both rates intact, as the $T_1$-driven
asymmetry of real readout demands. The lead example, the reliability of a
syndrome bit, is exactly $\tfrac12\bigl(1-\prod_i(1-2q_i)\bigr)$, a single product
that the field's ``$n$ times the per-measurement error'' rule overestimates by
tens of percent at modest sizes and rates, and that the asymmetric form
generalizes to inhomogeneous, relaxation-biased ancillas.

The calculus is deliberately bounded: forward rather than inverse,
complementary to matrix-inversion mitigation; valid under per-bit independence,
with correlated-crosstalk extension left open; and faithful only on the readout
sector. Its value is that it is \emph{reusable} and \emph{symbolic}, the same
object that governs Bloom filters and cipher maps now budgets the error of a
quantum-measurement post-processing pipeline, asymmetry preserved, without ever
forming the exponential assignment matrix. The natural next steps are a
correlated-channel (CTMP-regime) lift of the same operations and an empirical
calibration of predicted against measured derived-bit rates.

\paragraph{Acknowledgments.} This work is part of the Bernoulli type theory
program; the readout connection was developed in a concept note dated 2026-06-06.

\bibliography{references}

\end{document}
